\documentclass[11pt]{article}
\usepackage[margin=1in]{geometry}
\usepackage{amsmath,amssymb,amsthm}
\usepackage{xurl}
\usepackage{hyperref}
\sloppy
\let\oldus\_
\renewcommand{\_}{\oldus\allowbreak}

\newtheorem{theorem}{Theorem}
\newtheorem{lemma}[theorem]{Lemma}
\theoremstyle{remark}
\newtheorem{remark}[theorem]{Remark}
\newcommand{\R}{\mathbb{R}}
\newcommand{\cond}{\operatorname{cond}}

\title{Conjecture 00000005120:\\ same spectrum, exponentially different condition numbers\\
of the matrix inverse, separated by two Jordan structures}
\author{Jackmeson1 (AI-assisted)}
\date{2026-10-05}

\begin{document}
\maketitle

\section*{Statement}
\emph{English.} Definition: The spectrum and the function condition number are two layers
of sensitivity. Conjecture: There exist two matrices with the same spectrum but
exponentially different matrix-function condition numbers, and the separation is realized
by an explicit pair of Jordan structures.
(The Chinese text says the same: two matrices with the same spectrum whose matrix-function
condition numbers differ exponentially, the separation realized by an explicit pair of
Jordan structures.)

\textbf{Verdict: proved.} The statement is existential. ``Exponentially different'' needs a
growing parameter, so we give an explicit family indexed by the dimension $n\ge2$. Everything
below is the single Lean theorem \texttt{Conjecture5120.conjecture\_5120}, checked against
Mathlib.

\section{Definitions and reading}
\begin{itemize}
\item \emph{Matrices and spectrum.} Real $n\times n$ matrices. The spectrum is Mathlib's
\texttt{spectrum R M} (the $\mu$ with $\mu I-M$ not invertible). We also prove equality of
characteristic polynomials, so the eigenvalues agree with algebraic multiplicity (over
$\mathbb C$ as well, since the polynomial splits).
\item \emph{Jordan structures.} $J_n(\lambda)=\lambda I+N$, where $N$ has ones on the
superdiagonal (\texttt{jordanBlock}), and the scalar Jordan matrix $\lambda I$ with $n$ blocks
of size $1$ (\texttt{scalarJordan}).
\item \emph{Matrix function.} $f(x)=1/x$, whose matrix function is $f(X)=X^{-1}$ (Mathlib
\texttt{Ring.inverse}).
\item \emph{Matrix norm.} $\|M\|_\infty=\max_i\sum_j|M_{ij}|$, the operator norm induced by the
sup norm on $\R^n$ (Mathlib \texttt{Matrix.linftyOpNormedRing}; Lean type \texttt{Mat n}).
\item \emph{Condition numbers} (the standard definitions, as in Higham, \emph{Functions of
Matrices}, Ch.~3). The absolute condition number of $f$ at $X$ is the operator norm of the
Fr\'echet derivative, $\cond(f,X)=\|L_f(X)\|=\max_{Z\neq0}\|L_f(X)Z\|/\|Z\|$ (Lean
\texttt{cond}, defined as \texttt{fderiv}'s operator norm). The relative condition number is
$\cond_{\mathrm{rel}}(f,X)=\cond(f,X)\,\|X\|/\|f(X)\|$ (Lean \texttt{condRel}). Wikipedia
(``Condition number'') states the same: in several variables the condition number uses the
norm of the derivative, ``$J(x)$ denotes the Jacobian matrix of partial derivatives of $f$ at
$x$, and $\|J(x)\|$ is the induced norm on the matrix.''
\end{itemize}
The Fr\'echet derivative of inversion is Mathlib's \texttt{fderiv\_inverse}:
$L_{\mathrm{inv}}(X)Z=-X^{-1}ZX^{-1}$.

\section{Theorem}
For $n\ge1$ let $A_n=\tfrac12 I_n$ ($n$ Jordan blocks of size $1$) and
$B_n=J_n(\tfrac12)=\tfrac12 I_n+N_n$ (one Jordan block of size $n$).

\begin{theorem}[\texttt{conjecture\_5120}]\label{thm}
For every $n\ge2$:
\begin{enumerate}
\item $\chi_{A_n}=\chi_{B_n}=(X-\tfrac12)^n$ and $\sigma(A_n)=\sigma(B_n)=\{\tfrac12\}$;
\item $A_n$ and $B_n$ are not similar (different Jordan structures);
\item $\cond(\mathrm{inv},A_n)=4$ and $\cond(\mathrm{inv},B_n)\ge4^n$, so
$\cond(\mathrm{inv},B_n)\ge4^{\,n-1}\cond(\mathrm{inv},A_n)$;
\item $\cond_{\mathrm{rel}}(\mathrm{inv},A_n)=1$ and $\cond_{\mathrm{rel}}(\mathrm{inv},B_n)\ge2^n/4$.
\end{enumerate}
\end{theorem}

\section{Proof}
\textbf{(1) Same spectrum.} Both matrices are upper triangular with every diagonal entry
$\tfrac12$ (\texttt{A\_upper}, \texttt{B\_upper}), so Mathlib's
\texttt{charpoly\_of\_isUpperTriangular} gives $\chi=(X-\tfrac12)^n$ for both
(\texttt{charpoly\_A}, \texttt{charpoly\_B}). For a field, $\mu\in\sigma(M)$ iff $\mu$ is a root
of $\chi_M$ (\texttt{mem\_spectrum\_iff\_isRoot\_charpoly}); the only root of $(X-\tfrac12)^n$,
$n\ne0$, is $\tfrac12$ (\texttt{spectrum\_of\_charpoly}).

\textbf{(2) Different Jordan structures.} For $n\ge2$, $N_n\ne0$ (its $(0,1)$ entry is $1$).
If $B_n=PA_nP^{-1}$ then, since $A_n$ is scalar, $PA_nP^{-1}=A_n$, so $N_n=B_n-A_n=0$, a
contradiction (\texttt{not\_similar}).

\textbf{(3) Absolute condition numbers.} $A_n^{-1}=2I$, so $L(A_n)Z=-4Z$ and
$\|L(A_n)\|=4$: the bound $\|4Z\|\le4\|Z\|$ gives $\le4$, and any $Z\ne0$ gives $\ge4$
(\texttt{cond\_A}). For $B_n$ the inverse is explicit:
\[
(B_n^{-1})_{ij}=c(i,j)=\begin{cases}(-1)^{j-i}\,2^{\,j-i+1}, & i\le j,\\ 0,& i>j,\end{cases}
\]
proved by checking $B_nC=I$ entrywise (\texttt{key}, \texttt{B\_mul\_Binv}): row $i$ of $B_nC$ is
$\tfrac12c(i,j)+c(i+1,j)$, which is $1$ for $j=i$ and, for $j=i+e+1$,
$\tfrac12(-1)^{e+1}2^{e+2}+(-1)^e2^{e+1}=0$. Take $Z=E_{n-1,0}$ (a single $1$ in position
$(n-1,0)$), with $\|Z\|_\infty\le1$ (\texttt{norm\_single\_le}). Then
$(B_n^{-1}ZB_n^{-1})_{0,n-1}=c(0,n-1)^2=(2^n)^2=4^n$ (\texttt{entry\_calc}). Since every entry is
bounded by the $\infty$-norm (\texttt{entry\_le\_norm}),
$4^n\le\|L(B_n)Z\|\le\|L(B_n)\|\,\|Z\|\le\|L(B_n)\|$ (\texttt{cond\_B}).

\textbf{(4) Relative condition numbers.} $\|\tfrac12I\|=\tfrac12$ and $\|2I\|=2$
(\texttt{norm\_scalar}), so $\cond_{\mathrm{rel}}(\mathrm{inv},A_n)=4\cdot\tfrac12/2=1$
(\texttt{condRel\_A}). For $B_n$: $\|B_n\|\ge|(B_n)_{00}|=\tfrac12$ (\texttt{norm\_B\_ge}), and
$|c(i,j)|\le2^{j+1}$ gives every row sum of $B_n^{-1}$ at most
$\sum_{j<n}2^{j+1}=2^{n+1}-2$, so $0<\|B_n^{-1}\|\le2^{n+1}$ (\texttt{norm\_Binv\_le},
\texttt{norm\_Binv\_pos}). Hence $\cond_{\mathrm{rel}}(\mathrm{inv},B_n)\ge
4^n\cdot\tfrac12/2^{n+1}=2^n/4$ (\texttt{condRel\_B}).
\qed

\section{Scope and conventions}
\begin{itemize}
\item The matrix function is fixed as $f(x)=1/x$; the conjecture asks for some
matrix-function condition number, and inversion is the classical one. The exponential gap is
not claimed for every $f$. For $f=\exp$, for example, the same pair is \emph{not} separated
exponentially (remark, not formalized): $\|e^{sN_n}\|_\infty\le e^{s}$ for $0\le s\le1$, and the
integral formula $L_{\exp}(X)Z=\int_0^1e^{sX}Ze^{(1-s)X}\,ds$ then keeps
$\cond(\exp,B_n)\le e^{3/2}$ for every $n$.
\item ``Exponentially different'' is read as a ratio that grows exponentially in a parameter
of an explicit family. Here the parameter is the dimension $n$, and the ratios are at least
$4^{n-1}$ (absolute) and $2^{n}/4$ (relative). The eigenvalue $\tfrac12$ is fixed.
\item The matrix norm is the $\infty$-operator norm. Other norms change the constants. The
lower bound $\cond(\mathrm{inv},B_n)\ge4^n$ uses only one entry of $B_n^{-1}ZB_n^{-1}$ and
$\|E_{n-1,0}\|_\infty=1$. The Lean statement is proved for the $\infty$-norm only.
\item The spectrum is taken over $\R$; equal characteristic polynomials give equal complex
spectra and multiplicities as well.
\end{itemize}

\section{Lean formalization and axiom audit}
File \texttt{lean/Conjecture5120/Basic.lean} (Lean 4, Mathlib v4.33.1, only
\texttt{import Mathlib}; 403 lines). Main theorem:
\begin{verbatim}
theorem conjecture_5120 (n : Nat) (hn : 2 <= n) :
    A n = scalarJordan n (1/2) /\ B n = jordanBlock n (1/2) /\
    (A n).charpoly = (B n).charpoly /\
    spectrum R (A n) = {1/2} /\ spectrum R (B n) = {1/2} /\
    (not exists P : (Matrix (Fin n) (Fin n) R)^x, B n = P * A n * P^-1) /\
    cond n (inv n) (A n) = 4 /\ 4^n <= cond n (inv n) (B n) /\
    4^(n-1) * cond n (inv n) (A n) <= cond n (inv n) (B n) /\
    condRel n (inv n) (A n) = 1 /\ 2^n / 4 <= condRel n (inv n) (B n)
\end{verbatim}
Here \texttt{cond n f X := ||fderiv R f X||} (operator norm of the Fr\'echet derivative on
\texttt{Mat n}, the real $n\times n$ matrices with the $\infty$-operator norm),
\texttt{condRel n f X := cond n f X * ||X|| / ||f X||}, and \texttt{inv n := Ring.inverse}.
\texttt{lake build} succeeds, and \texttt{\#print axioms Conjecture5120.conjecture\_5120} reports
\texttt{[propext, Classical.choice, Quot.sound]}. There is no \texttt{sorry}, no
\texttt{native\_decide} and no added axiom.

\end{document}
